\documentclass[12pt,a4paper]{article}

\usepackage[utf8]{inputenc}
\usepackage[T1]{fontenc}
\usepackage[margin=2.5cm]{geometry}
\usepackage{hyperref}
\usepackage{microtype}
\usepackage{parskip}
\usepackage{lmodern}

\hypersetup{colorlinks=true, urlcolor=blue!60!black}

\pagestyle{empty}

\begin{document}

\begin{flushright}
Kundai Farai Sachikonye \\
Auf der Lichtung 19, 82178 Puchheim, Germany \\
\href{https://github.com/fullscreen-triangle}{github.com/fullscreen-triangle} \\
\texttt{kundai.sachikonye@bitspark.com} \\
(+49) 1626386344 \\[6pt]
\today
\end{flushright}

\vspace{1em}

Natera --- Hiring Team \\
Machine Learning Scientist, Multimodal AI \\
Remote \\
Job ID: 6004385004

\vspace{1.5em}

\textbf{Re: Machine Learning Scientist, Multimodal AI}

\vspace{1em}

Dear Hiring Team,

I am applying for the Machine Learning Scientist, Multimodal AI position. The
core of the role --- deep-learning models that integrate imaging, molecular, and
clinical data for oncology diagnostics and tumour-informed MRD --- is the area I
have spent the last several years building in, and the modalities you list map
closely onto systems I have already developed and deployed.

\textbf{Multimodal and molecular integration.}
This is my central focus. \href{https://github.com/fullscreen-triangle/gospel}{Gospel}
integrates host genotype, transcriptomics, and metagenomics in a joint
Bayesian-network model, validated against 1000~Genomes, gnomAD, and UK~Biobank.
\href{https://github.com/fullscreen-triangle/lavoisier}{Lavoisier} is a
production-grade mass-spectrometry pipeline whose CNN spectral-classification
stage (PyTorch) runs alongside a numerical pipeline under Ray-distributed
execution, reaching 98.9\% NIST accuracy on $>$100\,GB datasets.
\href{https://syndrome-seven.vercel.app/tools}{Syndrome} applies these methods to
oncology directly: neoantigen identification and MHC-binding prediction
(AUC~0.81 on 2{,}847 IEDB peptides) and tumour disease-state characterisation.
PyTorch and production-level Python are home ground.

\textbf{Imaging and foundation models.}
The imaging side connects through my work on microscopy image analysis ---
segmentation, deconvolution, and quantitative morphometry with metric-grounded
measurement --- which is directly adjacent to digital pathology and whole-slide
image analysis. On foundation models, I have fine-tuned pretrained models for
biomedical use through low-rank (LoRA) adaptation on physical observables, and I
use CNNs, attention, and transformer components across the pipelines above.

\textbf{Research to deployment.}
Every system I build is taken to deployment: publicly documented, validated
against reference data, and shipped as a browser-native or GPU-accelerated
application. Translating a prototype into a reproducible, validated tool is the
default of how I work, not an afterthought --- which is the part of this role I
would value most.

\textbf{Background.}
I hold an MSc in Computational Biology (Universit\"at T\"ubingen) and an MSc in
Pharmaceutical Biotechnology (HAW Hamburg), and I conducted doctoral research in
Computational Lipidomics at TUM (2019--2021) before moving to independent
research. English is native. I am based in Germany and work routinely in
international, English-language settings; I am glad to discuss working
arrangements and time-zone overlap.

I would welcome the opportunity to discuss how this multimodal, oncology-focused
background could contribute to Natera's diagnostic platforms.

\vspace{1.5em}
Yours sincerely,

\vspace{2.5em}
\textbf{Kundai Farai Sachikonye}

\end{document}
